\section{The hyperplane obstruction}\label{sec:hyperplane}

This section sets the constructions on abelian varieties aside and turns to
the other classical approach: induction on the dimension of the variety
through hyperplane sections. One places the variety in a family of hyperplane
sections, uses the algebraicity of the Hodge locus to control where a class
stays of type $(p,p)$, appeals to the conjecture in one dimension lower, and
transports the resulting cycles back up. Each ingredient of such
an argument is a theorem: the Lefschetz hyperplane theorem, hard Lefschetz,
the algebraicity of the Hodge locus \cite{CDK95}, and the properness of
relative Hilbert schemes. The question is whether they compose into an
induction. They do not, for an elementary reason that applies to every
argument of this shape, however the family of hypersurfaces is chosen.

Throughout, $X$ is a smooth complex projective variety of dimension $d$,
$j\colon Y\hookrightarrow X$ is a smooth hypersurface in the linear system of
an ample divisor, and $L=c_{1}(\cO_{X}(Y))\in H^{2}(X,\QQ)$.

\subsection{Two classical inputs}

\begin{theorem}[Lefschetz hyperplane theorem]\label{thm:LHT}
The restriction $j^{*}\colon H^{k}(X,\QQ)\to H^{k}(Y,\QQ)$ is an isomorphism
for $k\le d-2$ and is injective for $k=d-1$.
\end{theorem}

This is the topological form of the theorem. The standard proof, due to
Andreotti and Frankel, is by Morse theory on the smooth affine variety
$X\smallsetminus Y$; see \cite[Chapter~1]{Voi03}.

\begin{lemma}[Gysin and restriction]\label{lem:gysinrestrict}
For every $k$ and every $x\in H^{k}(X,\QQ)$,
\[
  j_{*}j^{*}x=x\cup L .
\]
Moreover $j_{*}\colon H^{k}(Y,\QQ)\to H^{k+2}(X,\QQ)$ is injective for
$k\ge d$.
\end{lemma}

\begin{proof}
The first identity is the projection formula together with
$j_{*}1=[Y]=L$. For the second, Poincar\'e duality on $X$ and on $Y$
identifies $j_{*}\colon H^{k}(Y)\to H^{k+2}(X)$ with the dual of
\[
  j^{*}\colon H^{2d-k-2}(X)\longrightarrow H^{2(d-1)-k}(Y)=H^{2d-k-2}(Y),
\]
so $j_{*}$ is injective on $H^{k}(Y)$ if and only if this $j^{*}$ is
surjective. By \Cref{thm:LHT} that $j^{*}$ is surjective whenever
$2d-k-2\le d-2$, that is whenever $k\ge d$.
\end{proof}

\subsection{The middle degree}

\begin{theorem}[Hyperplane obstruction]\label{thm:hyperplane}
Let $d=2p$ and let $\alpha\in H^{2p}(X,\QQ)$ be primitive. Then
$j^{*}\alpha=0$ for every smooth ample hypersurface $j\colon Y\hookrightarrow X$.
\end{theorem}

The primitive middle degree is where the Hodge conjecture is hard, and it is
where the restriction map returns nothing. \Cref{fig:ladder} draws the
Lefschetz decomposition of an abelian fourfold and the part of it that
restriction annihilates. The proof takes three lines.

\begin{proof}[Proof of \Cref{thm:hyperplane}]
Let $d=2p$ and let $\alpha\in H^{2p}(X,\QQ)$ be primitive. A class of degree
$k\le d$ is primitive when $L^{d-k+1}$ annihilates it, so in the middle degree
primitivity means $L\alpha=0$, that is,
$H^{d}_{\mathrm{prim}}(X,\QQ)=\ker\bigl(L\colon H^{d}(X,\QQ)\to H^{d+2}(X,\QQ)\bigr)$.
By \Cref{lem:gysinrestrict},
\[
  j_{*}\bigl(j^{*}\alpha\bigr)=\alpha\cup L=0 .
\]
The class $j^{*}\alpha$ lies in $H^{d}(Y,\QQ)$, on which $j_{*}$ is injective
by the case $k=d$ of \Cref{lem:gysinrestrict}. Hence $j^{*}\alpha=0$.
\end{proof}

\begin{corollary}[Blindness in the middle degree]\label{cor:blind}
Let $d=2p$ and let $\alpha,\beta\in H^{2p}(X,\QQ)$ be primitive Hodge classes.
Then $j^{*}\alpha=j^{*}\beta=0$ for every smooth ample hypersurface $Y$. In
particular no property of the restricted classes distinguishes $\alpha$ from
$\beta$, and no statement about Hodge classes on hypersurfaces of $X$, however
strong, carries any information about the primitive middle cohomology of $X$.
\end{corollary}

\begin{remark}[The size of what is lost]\label{rem:sizeofloss}
The primitive middle cohomology is not a small correction term. When $X$ is
a smooth hypersurface of degree $m$ in $\PP^{d+1}$, the dimension of
$H^{d}_{\mathrm{prim}}(X,\QQ)$ grows like $m^{d+1}$; and when $d=2p$, the
Hodge conjecture for $X$ in degree $2p$ is the statement that the Hodge
classes in $H^{d}_{\mathrm{prim}}(X,\QQ)$ are algebraic, since the remaining
Hodge classes of degree $2p$ are the multiples of $L^{p}$. By
\Cref{cor:blind}, all of this cohomology is invisible to hyperplane
restriction.
\end{remark}

\subsection{Below the middle degree}

Below the middle degree the restriction map is injective, and in degrees at
most $d-2$ it is an isomorphism, so a class can be transported downwards. The
difficulty moves to the return trip.

\begin{definition}[Lefschetz standard conjecture]\label{def:Bconj}
For a smooth projective variety $X$ of dimension $d$ with ample class $L$,
hard Lefschetz gives isomorphisms
$L^{d-k}\colon H^{k}(X,\QQ)\to H^{2d-k}(X,\QQ)$, and makes $H^{*}(X,\QQ)$
a representation of $\mathfrak{sl}_{2}$ in which $L$ raises the degree by
$2$ and $H$ acts on $H^{k}$ by $k-d$. Let $\Lambda$ denote the unique
endomorphism of degree $-2$ with $[L,\Lambda]=H$. The Lefschetz standard
conjecture $B(X)$ asserts that $\Lambda$ is induced by an algebraic cycle on
$X\times X$; this is equivalent to Kleiman's formulations \cite{Kle68,Kle94}.
\end{definition}

\begin{lemma}[$\Lambda$ undoes $L$ on primitive classes]\label{lem:LambdaL}
Let $\alpha\in H^{k}(X,\QQ)$ be primitive with $k\le d$. Then
$\Lambda L\alpha=(d-k)\,\alpha$.
\end{lemma}

\begin{proof}
This follows from the relation $[L,\Lambda]=H$, where $H$ acts on $H^{k}$ by
the scalar $k-d$, and from $\Lambda\alpha=0$, which holds because the
primitive classes of degree at most $d$ are the lowest weight vectors of the
$\mathfrak{sl}_{2}$-representation. Explicitly,
$\Lambda L\alpha=L\Lambda\alpha-[L,\Lambda]\alpha=0-(k-d)\alpha
=(d-k)\alpha$.
\end{proof}

\begin{theorem}[The return trip]\label{thm:returntrip}
Let $2p\le d-2$, so that $j^{*}\colon H^{2p}(X,\QQ)\to H^{2p}(Y,\QQ)$ is an
isomorphism. If $j^{*}\alpha$ is algebraic on $Y$, then $L\alpha$ is algebraic
on $X$. The implication ``$L\alpha$ algebraic $\Rightarrow$ $\alpha$
algebraic'', asked of all smooth projective $X$ and all Hodge classes
$\alpha$, is
equivalent to the Lefschetz standard conjecture.
\end{theorem}

\begin{proof}[Proof of \Cref{thm:returntrip}]
Let $2p\le d-2$, so $j^{*}\colon H^{2p}(X,\QQ)\to H^{2p}(Y,\QQ)$ is an
isomorphism by \Cref{thm:LHT}. If $j^{*}\alpha=\cl(W)$ for some
$W\in\CH^{p}(Y)_{\QQ}$, then $j_{*}W\in\CH^{p+1}(X)_{\QQ}$ and
\[
  \cl(j_{*}W)=j_{*}j^{*}\alpha=L\alpha
\]
by \Cref{lem:gysinrestrict} and the compatibility of the cycle class map with
proper push-forward, so $L\alpha$ is algebraic.

For the second statement, suppose $B(X)$ holds and $\alpha$ is primitive of
degree $2p<d$ with $L\alpha$ algebraic. Then $\Lambda$ is induced by an algebraic correspondence, hence
carries algebraic classes to algebraic classes, and
$\alpha=(d-2p)^{-1}\Lambda(L\alpha)$ by \Cref{lem:LambdaL}, so $\alpha$ is
algebraic. The same holds without primitivity. Write $P^{m}$ for the
primitive part of $H^{m}(X,\QQ)$. Iterating the relation $[L,\Lambda]=H$ as in
the proof of \Cref{lem:LambdaL} shows that $\Lambda L$ acts on the summand
$L^{i}P^{2p-2i}$ of the Lefschetz decomposition of $H^{2p}(X,\QQ)$ by the
scalar $(i+1)(d-2p+i)$, which is nonzero because $2p<d$. So $\Lambda L$ is
invertible on $H^{2p}(X,\QQ)$, its inverse there is a polynomial in
$\Lambda L$ with rational coefficients, and
$\alpha=(\Lambda L)^{-1}\Lambda(L\alpha)$ is algebraic.

Conversely, suppose that for every smooth projective variety and every Hodge
class $\alpha$ on it the algebraicity of $L\alpha$ implies that of $\alpha$.
Applied repeatedly, this shows that on every smooth projective variety of
dimension $e$ the hard Lefschetz isomorphisms
$L^{e-k}\colon H^{k}\to H^{2e-k}$, $k\le e$, carry the algebraic classes of
degree $k$ onto those of degree $2e-k$, since the preimage of an algebraic
class under $L^{e-k}$ is a Hodge class. This is Kleiman's conjecture $A$ for
every smooth projective variety, and $A$ for $X\times X$ implies $B(X)$
\cite{Kle68,Kle94}.
\end{proof}

Together the two theorems give the following.

\begin{corollary}[No dimensional induction through hyperplane sections]%
\label{cor:noinduction}
Let $X$ be smooth projective of dimension $d$ and let $\alpha\in\Hdg^{p}(X)$.
An argument that deduces the algebraicity of $\alpha$ from the Hodge
conjecture for smooth ample hypersurfaces of $X$, using only $j^{*}$ and
$j_{*}$, returns no information when $2p=d$ and $\alpha$ is primitive, and
otherwise establishes the algebraicity of $L\alpha$ rather than of $\alpha$.
\end{corollary}

\begin{proof}[Proof of \Cref{cor:noinduction}]
Suppose the argument uses only $j^{*}$ and $j_{*}$ to relate classes on $X$ to
classes on smooth ample hypersurfaces $Y\subset X$. First let $2p=d$ with
$\alpha$ primitive. Then $j^{*}\alpha=0$ by \Cref{thm:hyperplane}. A cycle on
$Y$ of codimension $p-1$ has its class in $H^{d-2}(Y,\QQ)$, which is
$j^{*}H^{d-2}(X,\QQ)$ by \Cref{thm:LHT}, so by \Cref{lem:gysinrestrict} its
image under $j_{*}$ lies in $L\cdot H^{d-2}(X,\QQ)$. By the Lefschetz
decomposition this subspace meets the primitive part of $H^{d}(X,\QQ)$ only
in $0$, so $j_{*}$ produces $\alpha$ only when $\alpha=0$, and no conclusion
about $\alpha$ is available. In the remaining cases the best that
$j_{*}j^{*}$ produces is $L\alpha$, and \Cref{thm:returntrip} identifies the
passage from $L\alpha$ to $\alpha$ with the Lefschetz standard conjecture.
\end{proof}

\subsection{Two gaps in inductive arguments}

In practice an induction of this shape breaks at one of two steps, each of
which looks harmless.

\begin{remark}[The dimension bookkeeping]\label{rem:dimensionbookkeeping}
A common formulation places $X$ in a Lefschetz pencil by choosing a generic
pencil $\{X_{t}\}_{t\in\PP^{1}}$ of hyperplane sections \emph{of $X$} and then
setting $X=X_{t_{0}}$ for some $t_{0}$. The two requirements are
incompatible: the members of a pencil of hyperplane sections of $X$ have
dimension $d-1$, so none of them is $X$. An induction that reads
$\dim X_{t}=d-1<d$ in its inductive step, and $X_{t_{0}}=X$ when it transports
the resulting cycle back, is using both readings of the same symbol.

Neither way of resolving the ambiguity rescues the argument. If $X$ is
required to be a member of the pencil, then the pencil lives in an ambient
variety of dimension $d+1$ and every member has dimension $d$; the inductive
parameter does not decrease and the induction is circular. If instead
$\dim X_{t}=d-1$, then the class $\alpha$ lives on $X$ and not on the members,
and the only available transports are $j^{*}$ and $j_{*}$. This is the
situation of \Cref{thm:hyperplane,thm:returntrip}, and \Cref{cor:noinduction}
applies.
\end{remark}

\begin{remark}[Pointwise existence is not a family]\label{rem:pointwise}
The second gap appears in arguments that produce cycles fibrewise and then
assemble them. Suppose an inductive hypothesis furnishes, for each $t$ in the
base $S$ of a family, a cycle $Z_{t}$ with a prescribed class. It does not
follow that the Hilbert polynomial $t\mapsto P_{Z_{t}}$ is constructible, or
semicontinuous, or takes finitely many values. Those are theorems about a flat
family, and the input here is a bare existence statement at each point with
no coherence between points. Without constructibility there is no section of
a relative Hilbert scheme over a dense open subset, and therefore nothing for
the valuative criterion of properness to extend. The step looks safe because
the relative Hilbert scheme of each fixed Hilbert polynomial is proper over
$S$; but properness extends a section that exists over a punctured curve, and
supplies no section where none is given.
\end{remark}

\begin{figure}[htbp]
\centering
  \includegraphics[width=\textwidth]{fig_ladder}
\caption{The Betti numbers $\binom{8}{k}$ of an abelian fourfold, drawn as
prisms and divided according to the Lefschetz decomposition
$H^{k}=\bigoplus_{i}L^{i}P^{k-2i}$: the primitive part $P^{k}$ is drawn dark,
and each image $L^{i}P^{k-2i}$ is drawn light in the colour of $P^{k-2i}$. The
arcs are the hard Lefschetz isomorphisms $L^{4-k}\colon H^{k}\to H^{8-k}$.
Below the middle degree the restriction $j^{*}$ to a smooth ample hypersurface
is injective and a class descends, but the return trip $j_{*}j^{*}$ multiplies
by $L$, and recovering the class is the Lefschetz standard conjecture; above
the middle degree a class is carried to the lower half by an arc. The
primitive middle cohomology $P^{4}$, of dimension $70-28=42$, is annihilated
by $j^{*}$ (\Cref{thm:hyperplane}), and neither an arc nor a restriction
reaches it.}
\label{fig:ladder}
\end{figure}
